% ============================================================
% Part III §III.1 — The Method of Reanalysis — DRAFT v1 (FCDP v2)
% Pack: plans/packs/part-iii-1-pack-v1.md (D1 embedded) · Quotes: part-iii-1-quotes.json
% Architecture written to Band G from the start (§III.0 lesson); M4 = Band R.
% ============================================================

\section{The Method of Reanalysis}

The method of this Part is the method of the dissertation carried to data, and its first word is a disclaimer of novelty: nothing in the frame is revised here, and nothing is added to it except an application layer---a set of stated jurisdictions by which each engine of the account meets each kind of evidence on terms fixed in advance. Four engines do the work. The analysis of boredom runs on the apparatus drawn from Heidegger's lecture course: the three forms of the attunement, their two structural moments, and a diagnostic grid that asks its questions in a fixed order---whether a determinate culprit can be named; whether an occupied surface stands over an emptied depth; whether the passing of time has ceased altogether while the while stretches. The analysis of design runs on the lens drawn from phenomenological design theory, which reads an environment as a set of conditions that afford or foreclose the chain's stages, and which will be held to its predictions when the deployment's complaints are counted. The game-analysis method developed for Part II's cases serves here in a deliberately bounded role, since its protocol wants a recorded episode of play---a trace in which an actualization can be followed beat by beat---and only the laboratory recordings supply one; where the data are testimony rather than trace, its vocabulary serves and its worksheet stays shut. Beneath all three, the metaphysics and the methodology of the earlier Parts govern the words: the chain from horizon to act, the union with its two poles of world-disclosedness and co-disclosedness, the economy of the image named by \textit{phantasia}-load. Jurisdiction, not doctrine, is what this section settles---which engine may speak of which data, and in what register of confidence.

That register is tiered, and the tiers are worn openly in the prose. Findings of the published studies are asserted as findings; the raw archive of the author's own laboratory work, which has not been published, is asserted with that status marked; and every reading that carries an empirical result over into the frame's categories is flagged as the frame's reading---a projection offered for its explanatory yield, never smuggled in as one more datum. The reader will therefore meet three voices in what follows: the record speaking, the record's unpublished margin speaking, and the frame speaking about both. Keeping them distinct is not decoration. It is what entitles the Part to its conclusions.

The founding constraint on all four engines comes from the lecture course itself, and it must be stated before any instrument is trusted, because it disciplines them all. Heidegger opens the analysis of boredom by refusing the very posture an empirical program most naturally assumes: ``«Q1»''\footnote{«Q1@».} An attunement is not an object lying present for inspection; it is the way a situation already has us, and the attempt to seize it as an object alters or destroys the thing seized---``«Q2»''\footnote{«Q2@».} Taken seriously, the constraint does not forbid measurement; it relocates what measurement measures. Electrodes, eye-trackers, and questionnaires never read the attunement itself; they read comportment---the fidgets and glances and answers in which an attunement casts its shadow---and the analysis that follows treats every signal accordingly, as shadow-evidence whose bearer must be reconstructed, never observed. A second consequence follows for classification. If boredom comes in forms rather than in quantities, then the forms are told apart by their structure and not by their intensity: the depth of an episode, the lecture course insists, ``«Q3»''\footnote{«Q3@».} A scale that runs from engaged to bored therefore begs the decisive question in advance, for it presumes that the phenomenon has one dimension, when the second form's whole character is to hold two dimensions apart---a surface that remains occupied and a depth that has been left empty. Hence the rule on which this Part's deepest claims will stand or fall: the second form of boredom is detectable only as the disagreement of a surface channel and a depth channel, never from either alone, and never from an instrument built to collapse the two into one score. The rule carries a corollary for testimony that must be honored even when it is inconvenient: the subject of the second form sincerely denies being bored, reading the occupied surface as the whole story---so self-report is a channel to be listened to, and never a verdict to be accepted.

A caution of category completes the constraint, and it is constitutive of the method rather than an apology appended to it. A seventeen-minute video does not disclose beings as a whole; a course module does not visit the profound boredom in which a world hangs suspended; and nothing in this Part pretends otherwise. What transfers from the lecture course to the laboratory is the structural grammar of the forms---culprit-structure, emptiness-structure, temporal structure---taken as a heuristic that must earn its keep case by case. The first two forms transfer cleanly, since their grammar is episodic on its face. The third transfers only as an analogy, explicitly marked wherever it appears, for a collapsed gaze may resemble the stance of profound boredom without any claim that a participant underwent a metaphysical disclosure between one film and the next.

With the constraint and the caution in place, the corpus's channels can be sorted in advance, so that no instrument acquires authority by accident. As surface channels the Part admits the gaze metrics, the recorded comportment, and every self-report of engagement or presence, closed or open---all of them readings of the occupied surface. As proxies for depth it admits three, unequal in strength: the electroencephalographic profile, which enters with corroborating weight only; the felt-duration reports, which enter as the strongest depth evidence the corpus owns, since the stretching of the while is the one signature that the second form cannot suppress; and, weaker still, the direct self-ratings of boredom and of hollowness---the closed rating where an instrument asks about boredom separately from engagement, and the hollowness that leaks through open testimony---admitted as depth traces only where they stand against a surface-positive answer from the same voice, and always under the denial-caution stated above, since the second form's subject reads the occupied surface as the whole story. Divergence---the method's true object---is admitted wherever a surface and a depth speak of the same episode: instrument against instrument in the laboratory, item against item within a single laboratory questionnaire, and, in the deployment testimony, only within a single student's answers, where praise and hollowness share one sheet. And some things are admitted as nothing at all: the content-quiz scores, which sat at ceiling and so carry no signal; the platform's automated item statistics, which assign correct answers to survey questions and are artifacts of the export rather than psychometrics; and any composite that would average engagement and boredom into a single number, which the rule above disqualifies in principle. The sorting is itself falsifiable, and should be seen to be: nothing guarantees that surface and depth will ever disagree, and had every channel of this corpus agreed with every other, the second form would simply have found no purchase here. The frame's design lens, likewise, is made to stake predictions---which frictions a student population would name---that the deployment section counts against the record. A frame that could not have failed would explain nothing by succeeding.

The deployment testimony, finally, was prepared for this method by a coding protocol whose reliability is stated rather than assumed. The instrument is a fourteen-item course questionnaire, stable across the three terms studied. Three content items sat at ceiling and were excluded, and one item is a survey preamble carrying no data. Two closed items carry the experience data: whether the student was able to use the platform, and whether it produced a feeling of immersion or presence, with a video-only answer available. Four open items invited prose; four demographic items close the form. Two hundred forty-one students answered across the terms, and their eight hundred seventy-three open responses form the coded corpus. The codebook holds eight families: boredom-form signatures, with sub-tags for the named frictions; flight patterns; temporal texture; design conditions keyed to the chain's stages; involvement-channel vocabulary; incorporation-positive signatures; suggestions; and an open family for emergent codes. Every response was coded at least once. A subsample of one hundred forty-three responses was coded twice, blind, and the two passes agreed at a mean Jaccard of .878 over full code sets; the four disagreements that survived rule-governed merging were adjudicated by hand and logged. Names were replaced by codes throughout, and references to instructors were masked. The record's limits are equally part of the record. The survey is one self-report channel, so the second form can appear in it only as an echo---a divergence-shaped pattern inside a single voice---and never as a detection. The questionnaire is retrospective, so the temporal texture of episodes is largely mute in it, and the laboratory archive must carry that thread. No outcome measure survives the ceiling, so nothing here speaks of learning gains. And the terms differ in composition, so their comparison stays directional. Within those limits, the corpus is ready; what remains is to read it.

% ============================================================
% FOOTER (FCDP provenance)
% Material: ALL NEW (greenfield; P7 = N/A). Retained/corrections: n/a. Omissions: none.
% Quotes: Q1–Q3 all used, footnote-wrapped at write time (§III.0 lesson).
% Counterargument dispositions: C1 ANSWER (category caution ¶4); C2 ANSWER (coding
%   record ¶6); C3 ANSWER (falsifiability, ¶5).
% GAUNTLET REPORT (final): G-A PASS-blended (avg 28.9, long .471; short .333
%   explained by Band-R M4 blend; nominalization 4.48; §III.0 waivers apply) ·
%   G-B PASS (3/3, exit 0) · G-C PASS (all loci from bank) · G-D PASS (greps
%   clean; zero pronoun tokens) · G-E PASS after 1 cycle (pack v1.1: merge-
%   procedure clause, artifact mechanism, preamble item; prose: three depth
%   proxies restored, 14th item accounted) · G-F attested · G-G PASS after
%   1 cycle ("a single student's answers" noun fix; 13-vs-14 cured by preamble).
%   R-cycles used: 1 of 3.
% ============================================================
